\chapter{中国科学院大学
本科生毕业论文（设计）撰写规范指导意见（节选）}{
% {
% \let\cleardoublepage\relax
% }
本科生的毕业论文（设计）是培养学生综合运用所学基本理论、专业知识、基本技能进行科学研究的重要环节，是本科学生毕业和授予学士学位的必要条件，也是衡量学校教育质量的重要内容。为保证本科生毕业论文（设计）的质量，根据《学位论文编写规则》（GB/T 7713.1—2006）、《信息与文献 参考文献著录规则》（GB/T 7714—2015）和《学术出版规范 期刊学术不端行为界定》（CY/T 174—2019）等国家有关标准，结合《中国科学院大学本科毕业论文（设计）管理规定》，特制定本指导意见。

\section{组成及要求}
毕业论文（设计）一般由以下几个部分组成：封面、原创性声明及授权使用声明、摘要、目录、符号说明（若有）、正文、参考文献、附录（若有）、致谢等。
\subsection{封面}
一律采用中国科学院大学规定的统一中英文封面，封面包含内容如下：

\begin{enumerate}
    \item 论文题目，应简明扼要地概括和反映整个论文的核心内容，一般不宜超过25个汉字（符），英文题目一般不应超过150个字母，必要时可加副标题。题目中应尽量避免使用缩略词、首字母缩写词、字符、代号和公式等。
    \item 作者姓名,根据《中国人名汉语拼音字母拼写规则》（GB/T 28039-2011），英文封面中的姓和名分写，姓在前，名在后，姓名之间用空格分开。姓和名需写全拼，姓全大写，名首字母大写。外国留学生姓名书写顺序以护照格式为准，字母全部大写。
    \item 指导教师，需同时填写导师姓名、专业技术职务和工作单位。如果有多位导师（均需经院系批准），第一导师在前，第二导师等依次在后。
    \item 学位类别，包括学科门类及学位级别。学科门类如理学、工学等。学位级别为学士。
    \item 专业，填写攻读专业全称，须与学籍信息一致，不可用简写。
    \item 学院（系），填写就读学院（系）全称，如中国科学院大学XX学院。
    \item 时间，填写论文提交学位授予单位的年月，使用阿拉伯数字标注。一般学士学位的申请安排在夏季，夏季申请学位的论文标注6月。例如：2023年6月。
\end{enumerate}

\subsection{原创性声明及授权使用声明}
本部分内容提供统一的模版，具体内容见样张3，提交时作者和导师须亲笔签名。
\subsection{摘要和关键词}
论文摘要包括中文摘要和英文摘要（Abstract）两部分。论文摘要应概括地反映出本论文的主要内容，说明本论文的主要研究目的、内容、方法、成果和结论。不宜使用公式、图表、表格或其他插图材料，不标注引用文献。英文摘要与中文摘要内容应保持一致。中英文摘要以300-500字为宜。

摘要最后注明本文的关键词（3～5个）。关键词是为了文献标引和检索工作，从论文中选取出来，用以表示全文主题内容信息的单词或术语。关键词以显著的字符另起一行并隔行排列于摘要下方，左顶格，中文关键词间用中文逗号隔开。英文关键词应与中文关键词对应，首字母应大写，用英文逗号隔开。

摘要应另起一页，与正文前的内容连续编页（用罗马字符）。
\subsection{目录}
目录应包括论文正文中的全部内容的标题，以及参考文献、附录（若有）和致谢等，不包括中英文摘要。目录页由论文的章、条、附录等序号、名称和页码组成。正文章节题名建议最多编到第三级标题，即×.×.×（如1.1.1）。一级标题顶格书写，二级标题缩进一个汉字符位置，三级标题缩进两个汉字符位置。论文中若有图表，应有图表目录，置于目录页之后，另页编排。图表目录应有序号、图题或表题和页码。

目录应另起一页，与正文前的内容连续编页（用罗马字符）。
\subsection{符号说明（若有）}
如果论文中使用了大量的物理量符号、标志、缩略词、专门计量单位、自定义名词和术语等，应编写成注释说明汇集表。若上述符号等使用数量不多，可以不设此部分，但必须在论文中首次出现时加以说明。
论文中若有符号说明，应置于目录之后、正文之前，另起一页，与正文前的内容连续编页（用罗马字符）。
\subsection{正文}
正文一般包括绪论、论文主体、研究结论与展望等部分。

提供以下两种结构供学生参考使用，例如：

第一种结构：第1章 绪论、第2章 材料与方法、第3章 结果、第4章 讨论、第5章 研究结论与展望；

第二种结构（适用于几部分相对独立又有联系的研究内容）：第1章 绪论、第2章 ××××（第一部分研究内容）、第3章 ××××（第二部分研究内容）、第4章 ××××（第三部分研究内容）、第5章 研究结论与展望。
\begin{enumerate}
    \item 绪论
    
    绪论应包括选题的背景和意义，国内外相关研究成果与进展述评，本论文所要解决的科学与技术问题、所运用的主要理论和方法、基本思路和论文结构等。绪论应独立成章，用足够的文字叙述，不与摘要雷同。要实事求是，不夸大也不弱化前人的工作和自己的工作。
    \item 论文主体
    
    论文主体是正文的核心部分，占主要篇幅，它是将学习、研究和调查过程中筛选、观察和测试所获得的材料，经过加工整理和分析研究，进而形成论点。由于不同学科专业及具体选题的差异，此部分不作统一规定。但总体内容必须实事求是，客观真切，准确完备，合乎逻辑，层次分明，简练可读。
    
    \item 研究结论与展望
    
    研究结论是对整个论文主要成果的总结，不是正文中各章小结的简单重复，应准确、完整、明确、精炼。包括对整个研究工作进行归纳和综合而得出的总结，还应包括所得结果与已有结果的比较和本课题尚存在的问题，以及进一步开展研究的见解与建议等。
\end{enumerate}
\subsection{参考文献}
本着严谨求实的科学态度撰写论文，凡学位论文中有引用或参考、借鉴他人思想或成果之处，均应按一定的引用规范，列于文末（通篇正文之后），参考文献部分应与正文的文献引用一一对应，注重合理引用，严禁抄袭剽窃等学术不端行为。
\subsection{附录（若有）}
主要列入正文内过分冗长的公式推导、供查读方便所需的辅助性数学工具或表格、数据图表、程序全文及说明、调查问卷、实验说明等。
\subsection{致谢}
对给予各类资助、指导和协助完成研究工作，以及提供各种对论文工作有利条件的单位及个人表示感谢。致谢应实事求是，切忌浮夸与庸俗之词。

致谢末尾应具日期，日期与论文封面一致。


\section{撰写要求}

\subsection{毕业论文的基本要求}
毕业论文必须是一篇系统的、完整的学术论文，遵循既定的学术规范与要求，不仅要符合学位论文的形式规范，更要符合学位论文的质量规范。做到：学术观点明确，立论正确，方法科学，材料翔实，数据可靠，推理严谨，论证充分，引用规范，结构合理，层次分明，文字通顺，表达准确，学风严谨。毕业论文应为学位申请者在导师的指导下独立完成的科学研究成果，没有伪造、篡改、剽窃、他人代写、论文买卖及其它学术不端行为。


\subsection{文字、标点符号和数字}

除特殊需要外，毕业论文（设计）一律用国家正式公布实施的简化汉字书写。标点符号的用法以《标点符号用法》（GB/T 15834—2011）为准。数字用法以《出版物上数字用法》（GB/T 15835—2011）为准。

为了便于国际合作与交流，毕业论文（设计）亦可有英文或其他文字的副本。

\subsection{论文正文}

\subsubsection{章节及各章标题}
论文正文须由另页右页（奇数页）开始，用阿拉伯数字连续编码，一直到全文最后。正文内部新章节无须另页右页（奇数页）开始。

若按照“绪论-材料与方法-结果-讨论-研究结论与展望”的形式撰写论文时，绪论、材料与方法、结果、讨论、研究结论与展望分别单独成为章节并作为章节标题使用。

各章标题中尽量不采用英文缩写词，对必须采用者，应使用本行业的通用缩写词。标题中尽量不使用标点符号。

\subsubsection{序号}
\begin{enumerate}
    \item 标题序号

论文标题分层设序。层次以少为宜，根据实际需要选择。各层次标题一律用阿拉伯数字连续编号。以三级标题为宜，最多四级。若确需要再增加一级，以小括号形式表示；不同层次的数字之间用小圆点“.”相隔，末位数字后面不加点号，如“1.1”，“1.1.1”等；章的标题居中排版，各层次的序号均左起顶格排，序号与题名间空一个汉字符。

    \item 图表等编号

论文中的图、表、附注、公式、算式等，一律用阿拉伯数字分章依序连续编码。其标注形式应便于互相区别，如：图1-1（第1章第一个图）、图2-2（第2章第二个图）；表3-2（第3章第二个表）等。附录的图表参考正文的编号方式，如附图1-1或附表1-1。

    \item 页码

正文页码从绪论开始按阿拉伯数字（1，2，3……）连续编排，页码应位居左页左下角、右页右下角；正文前的部分（中英文摘要、目录等）用大写罗马数字（I，II，III…）单独编排，页码位于页面下方居中。
\end{enumerate}
\subsubsection{页眉}
页眉从摘要开始，奇数页上标明“摘要”、“Abstract”、“目录”、“图表目录”等，偶数页上标明论文题目（英文摘要标明英文题目）。正文（即第1章开始到最后一章）的页眉，奇数页上标明每一章名称，偶数页上标明论文题目。参考文献、附录、致谢等的页眉，奇数页标明“参考文献”、“附录”、“致谢”等，偶数页标明论文题目。页眉居中设置。

\subsubsection{名词和术语}
科技名词术语及设备、元件的名称，应采用全国科学技术名词审定委员会公布的权威标准或其他相关权威信息源规定的术语或名称。标准中未规定的术语要采用行业通用术语或名称。全文名词术语必须统一。一些特殊名词或新名词应在适当位置加以说明或注解。双名法的生物学名部分均为拉丁文，并为斜体字。

采用英语缩写词时，除本行业广泛应用的通用缩写词外，文中第一次出现的缩写词应该用括号注明英文原词。新的外来名词应用括号注明英语全称和缩写语。

\subsubsection{量和单位}

量和单位要严格执行《国际单位制及其应用》（GB 3100-93）、《有关量、单位和符号的一般原则》（GB3101—93）有关量和单位的规定。量的符号一般为单个拉丁字母或希腊字母，并一律采用斜体（pH例外）。

\subsubsection{图和表}

论文中若有图和表，应设置图表目录，先列图后列表，置于目录页后，另页编排。

{(1) 图}

图片大小适当，图边界在页面范围内（图边界离页面边界距离大于页边距）。若图片中包含文字，文字大小不超过正文文字大小。

图包括曲线图、构造图、示意图、框图、流程图、记录图、地图、照片等，宜插入正文适当位置。引用的图必须注明来源。具体要求如下：
\begin{itemize}
    \item 图应具有“自明性”，即只看图、图题和图注，不阅读正文，就可理解图意。每一图应有简短确切的图题，连同图序置于图下居中。
    \item 图中的符号标记、代码及实验条件等，可用最简练的文字横排于图框内或图框外的某一部位作为图注说明，全文统一。图题建议用中文及英文两种文字表达。
    \item 照片图要求主要显示部分的轮廓鲜明，便于制版，如用放大、缩小的复制品，必须清晰，反差适中，照片上应有表示目的物尺寸的标尺。
    \item 图片一般设为高6cm×宽8cm，但高、宽也可根据图片量及排版需要按比例缩放。中文（宋体）英文（Times New Roman）图注为五号字，1.25倍行距。
    \item 文中尽量不用世界地图、全国地图。如果一定要用，凡涉国界图件（国内部分地区、全国、世界部分地区、全球）必须使用自然资源部标准地图底图（下载网址：http://bzdt.ch.mnr.gov.cn），所用底图边界要完全无修改（包括南海诸岛位置），为适应排版时图的缩放，比例尺一律用线段比例尺，而不用数字比例尺。并在图题下注明“注：该图基于自然资源部标准地图服务网站下载的审图号为GS(2021)××××号的标准地图制作，底图边界无修改。”
\end{itemize}


{(2) 表}

表的编排一般是内容和测试项目由左至右横读，数据依序竖排，应有自明性，引用的表必须注明来源。具体要求如下：
\begin{itemize}
    \item 每一表应有简短确切的题名，连同表序置于表上居中。必要时，应将表中的符号、标记、代码及需说明的事项，以最简练的文字横排于表下作为表注。表题建议用中文及英文两种文字表达。
    \item 表内同一栏数字必须上下对齐。表内不应用“同上”、“同左”等类似词及“″”符号，一律填入具体数字或文字，表内“空白”代表无此项，“—”或“…”（因“—”可能与代表阴性反应相混）代表未发现，“0”该表实测结果为零。表内未测出值可以用“N.D. ”表示。
    \item 表格尽量用“三线表”，避免出现竖线，避免使用过大的表格，确有必要时可采用卧排表，正确方位应为“顶左底右”，即表顶朝左，表底朝右。表格太大需要转页时，需要在续表表头上方注明“续表”，表头也应重复排出。
    \item 中文（宋体）英文（Times New Roman）表注为五号字，1.25倍行距。
\end{itemize}

\subsubsection{表达式}
论文中的表达式需另行起，原则上应居中。若有两个以上的表达式，应从“1”开始的阿拉伯数字进行编号，并将编号置于括号内。编号采用右端对齐。表达式较多时可分章编号。

较长的表达式如必须转行，只能在+，-，×，÷，＜，＞等运算符之后转行，序号编于最后一行右顶格。

\subsection{参考文献}
参考文献格式规范参照《信息与文献 参考文献著录规则》（GB/T 7714—2015），或可参照国际刊物通行的参考文献格式。各学科群分会可根据本学科的一般规范制定相应的参考文献格式。

文后参考文献和参考文献在正文中的标注方式可采用“顺序编码制”或“著者—出版年制”。确定采用某种方法后，文后参考文献和参考文献在正文中的标注方式要对应。

文后参考文献按“顺序编码制”组织时，各篇文献应按正文部分首次引用时标注的序号依次列出；文后参考文献按“著者—出版年制”组织时，条目不排序号，先按语种分类排列，语种顺序是：中文、日文、西文、俄文、其他文种；然后按著者字序和出版年排列。中文和日文按第一著者的姓氏笔画排序，中文也可按汉语拼音字母顺序排列，西文和俄文按第一著者姓氏字母顺序排列。当一个著者有多篇文献并为第一著者时，该著者单独署名的文献排在前面（并按出版年份的先后排列），接着排该著者与其他人合写的文献。

文后参考文献加标题“参考文献”，并列入全文目录。

凡正文里标注了参考文献的，其文献都必须列入文后参考文献。文后参考文献应集中著录于正文之后，不分章节著录。

正文中未被引用但被阅读或具有补充信息的文献可集中列入附录中，其标题为“荐读书目”。

详细内容请参考《中国科学院大学本科生毕业论文（设计）撰写规范指导意见》。

\section{排版与印刷要求}

\subsection{纸张与页面要求}
\begin{table}[h]
    \centering
        \bicaption{\enspace 排版和印刷要求}{\enspace Typography and Printing Requirements}
    \begin{tabular}{lc}
        \hline
        %\multicolumn{num_of_cols_to_merge}{alignment}{contents} \\
        %\cline{i-j}% partial hline from column i to column j
        项目名称 & 要求\\
        \hline
        纸张&A4（210mm×297mm），幅面白色\\
        页面设置&上、下2.54cm，左、右3.17cm，页眉、页脚距页边界1.5cm\\
        封面&采用国科大统一格式\\
        页眉&宋体小五号居中，英文和阿拉伯数字用Times New Roman体\\
        页码&Times New Roman体小五号 \\

        \hline
    \end{tabular}

    \label{tab:printrequirements}
\end{table}

其他要求详见《中国科学院大学本科生毕业论文（设计）撰写规范指导意见》。}